\documentclass[11pt]{article}
% Journal review draft. Legacy filename retained for build compatibility.
% This is not a claim of compliance with the publisher's current template.
\usepackage[T1]{fontenc}
\usepackage[utf8]{inputenc}
\usepackage[margin=1in]{geometry}
\usepackage{times,microtype,graphicx,array,amsmath,booktabs,longtable}
\usepackage[round,authoryear]{natbib}
\usepackage{xurl}
\usepackage[hidelinks]{hyperref}
\setlength{\emergencystretch}{2em}
\newcommand{\framework}{\textsc{dsh}}
\input{../generated/benchmark-metadata.tex}
\title{Evaluating Agent Skills for Version-Specific Plugin Migration: A Retrospective Study}
\author{Author information supplied separately}
\date{Review draft: 17 September 2026}
\begin{document}
\maketitle
\begin{abstract}
Agent skills package maintenance knowledge into reusable instructions and tools, but higher benchmark scores do not necessarily establish reliable migration behavior. We study a version-pinned plugin-upgrade skill through five historical configurations and a more completely archived, within-configuration comparison of 16 static tasks with two attempts per condition. Historical mean paired changes range from $-3.07$ to $+10.67$ points on recorded 0--100 reward scales; differences in task pools, harnesses, and graders prevent interpreting this range as a model capability curve. In the 64-output archived comparison, the skill increases mean recorded reward from 93.83 to 98.75, a gain of 4.92 points with a task-bootstrap interval of [0.31, 10.86]. Half the task pairs are at the ceiling in both conditions, and one task contributes disproportionately. A complete contract-stratified analysis accounts for all 328 original criterion decisions. A bounded, non-blind AI-assisted review of 10 answers and 56 criteria identifies three scoring disagreements affecting two answers; the authors subsequently reported human checking with broadly similar assessments. Preserving original scores and reporting alternative replacements yields intervals that touch or cross zero. Output inspection reveals errors involving public API boundaries, path containment, and process liveness that plausible migration advice can conceal. The study contributes a traceable maintenance case study and separates observed score improvements from operational correctness, resource use, and measurement uncertainty. It does not establish an inverted-U relationship or generalize effects across models.
\end{abstract}
\noindent\textbf{Keywords:} software maintenance; plugin migration; agent skills; retrospective evaluation; LLM-as-a-judge

\section{Introduction}
Software migration requires more than replacing names in source code. A repair must respect the contract of the target version, distinguish public APIs from internal implementation, and preserve runtime behavior. Plugin systems make these requirements particularly visible: a change can affect loading, authentication, rendering, event dispatch, or process shutdown independently. Advice that sounds consistent with a release note may still recommend an unavailable API or omit a boundary condition.

An agent skill bundles instructions, reference material, and sometimes executable tools for repeated use. Such a package can reduce the effort needed to locate version-specific facts. It can also introduce stale instructions, unnecessary work, or confidence in an incomplete repair. Existing skill benchmarks already demonstrate heterogeneous benefits and overhead \citep{skillsbench,swe-skills-bench,webdev-skills-bench}. Consequently, the contribution of a maintenance case study cannot rest only on showing that a skill sometimes helps.

We investigate the shipped plugin-upgrade skill for \framework{}, focusing on what its recorded improvements actually support. The available evidence accumulated during development rather than under one prospectively controlled experiment. We retain this history, distinguish incompatible protocols, and give greatest analytical attention to a later comparison with archived outputs, per-criterion judgments, and execution records. The study asks:
\begin{enumerate}
\item RQ1: What within-configuration changes in recorded task reward accompany skill availability?
\item RQ2: What migration-contract failures and grading disagreements appear in a bounded review of recorded answers?
\item RQ3: How do resource accounting and grading sensitivity qualify the apparent benefit?
\end{enumerate}

Our contributions are a version-sensitive migration evaluation with explicit task and exposure provenance; a complete contract-stratified account of the 64-output focal comparison, separated from heterogeneous historical context; and grounded case analysis connecting apparent success to API, containment, and liveness requirements. The review provides counterexamples to overly broad correctness claims, not an estimate of failure prevalence. We do not claim a new general law relating model capability to skill benefit.

\section{Related Work}\label{sec:related}
\paragraph{Skills and software engineering.}
SkillsBench evaluates skills across diverse tasks \citep{skillsbench}. SWE-Skills-Bench examines their utility in software engineering \citep{swe-skills-bench}. SkillLens studies skill generation and consumption \citep{skilllens}, while WebDev-Skills-Bench compares skill conditions in web development, including controls for supplied context \citep{webdev-skills-bench}. These studies make heterogeneous effectiveness and resource overhead established questions. Our narrower contribution is an artifact-traceable account of version-contract errors and evaluation sensitivity in one maintained plugin-migration setting. We do not replicate their breadth or their controlled ablations.

\paragraph{Version-sensitive maintenance.}
VersiCode studies version-controllable code generation, and CODEMENV studies code migration \citep{versicode,codemenv}. General software-agent evaluations such as SWE-bench and AgentBench emphasize task and environment design \citep{swe-bench,agentbench}. Our setting combines migration diagnosis with historical executable tasks, but the most completely archived comparison consists of static answers. Its reward therefore cannot be treated as an executable repair success rate.

\paragraph{Measurement validity.}
LLM judges enable scalable evaluation but require validation \citep{mt-bench-judge}; self-preference is a documented concern \citep{self-preference}. This matters when solvers and judges share a model family. Development exposure is a separate issue: repeated use of the same task during skill refinement does not create a fresh holdout \citep{time-travel-contamination}. Our audit addresses selected observable errors and provenance. It cannot eliminate grader bias, task-selection bias, or exposure effects.

\section{Setting and Evidence}
\subsection{Versioned migration and the intervention}
\framework{} is a plugin-based framework with a marketplace and frequent releases. Relevant changes include renamed exports, authentication boundaries, build-artifact layouts, session representations, and version-routing rules. Tasks specify an initial state, a target contract, permitted operations, and a deliverable. Static tasks request diagnosis; hands-on tasks require a repair in a version-pinned environment. A host cold boot checks liveness, while task-specific probes exercise additional requirements. Neither establishes correctness beyond the checked behaviors.

Upgrade cards record symptoms, causes, actions, version applicability, and provenance. For example, card \texttt{DSH-0.1.2-A1-01} concerns removal of the Host APIProxy surface on the 0.1.1-rc.2 to 0.1.2-alpha.1 edge. Correct migration requires locating actual call sites and selecting the corresponding Remote interfaces, rather than applying a client-plane recipe to every plugin. This illustrates why merely mentioning the correct release or card is insufficient.

The shipped skill links inspection, version-corridor analysis, implementation, and verification. Cards can inform both the skill and the benchmark. Thus the comparison is an open-book evaluation of the complete package, including access to facts, rather than an isolated test of procedural reasoning. Historical repository states differ, so the intervention is not assumed byte-identical across configurations.

\subsection{Task construction and exposure}
Contributors turn migration incidents and version changes into cards, fixtures, instructions, reference solutions, and graders. Contributor self-check and maintainer review constitute development quality control, not independent blinded annotation. The repository's exposure ledger distinguishes feedback used to refine the skill, grader development, and uncleared exposure. We make no claim that the evaluated tasks form an independent holdout.

An immutable earlier snapshot contains \BenchmarkTaskCount{} tasks (\BenchmarkStaticCount{} static and \BenchmarkHandsOnCount{} hands-on). This inventory is not the denominator of every experiment. Historical configurations use 56, 23, 21, or 22 scored tasks. The focal comparison uses a recorded selection of 16 static tasks. Repeated attempts are nested within tasks; they do not increase the number of independent migration problems. Related tasks may also share incident sources, limiting task-level uncertainty estimates.

\subsection{Evidence tiers}
We distinguish three tiers by what the archive supports.
\begin{enumerate}
\item Historical context: five within-configuration paired summaries, with heterogeneous task pools, aggregation, grading, and raw-artifact coverage.
\item Focal archived comparison: 16 tasks, two conditions, and two repeats, with 64 reports, original verdicts, scores, selection and execution records, and an offline integrity audit.
\item Supplementary evidence: a newer 16-task Qwen comparison with scored records but without archived original answers and judgment reasons, plus other historical runs that are not pooled into the focal estimate.
\end{enumerate}
A complete set of files improves auditability; it does not make a retrospective selection prospective or establish the validity of its grader.

\section{Method}
\subsection{Historical paired summaries}\label{sec:paired}
For each historical configuration, we use its recorded task reward on a 0--100 scale and preserve its original within-task aggregation. Qwen uses means of three scored attempts, DeepSeek and the two GLM groups use medians of three runs or rounds, and Terra uses one attempt. These numeric scales do not imply equivalent constructs: keyword grading, semantic judgment, and composite executable rewards differ. We do not pool configurations or regress their gains on their baseline scores as an independent capability measure.

Historical percentile intervals use 10,000 task-paired bootstrap samples with mulberry32 seed 20260907. The reported two-sided Wilcoxon signed-rank calculations exclude zero differences and use a tie-corrected normal approximation with continuity correction. These exploratory summaries can disagree near a boundary; neither is a substitute for validating the endpoint. No multiple-comparison-adjusted confirmatory family is claimed.

\subsection{Focal comparison and estimand}
The focal run is the archived GLM-5.3-Flash unified S16 comparison. The model label is the recorded gateway label, not an independently verified served-weight identity. Each task has two no-skill and two with-skill outputs, run through the ZCode subagent host. The conditions share a prompt template, but workspace paths differ. This is not a container-isolated or byte-identical-prompt experiment. The selection script draws 16 of S1--S22 using seed 20260915, stratifying by problem family and fixture-size tercile with proportional quotas. Balance constraints require at least two hard tasks, two Chinese-prompt tasks, and two large-fixture tasks. The selector reads metadata and fixture sizes rather than scores; it excludes S3, S4, S7, S8, S10, and S16. The selected set contains five static-contract diagnosis, seven runtime/client API, three release/install, and one profile/Cordis task; 14 prompts are English and two are Chinese. This reproducible draw is within a development-exposed pool, not a random sample of migration problems in general. The schedule alternates arm order in task blocks and reverses it in round two; the recorded execution order is retained. Four pilot cells on excluded S4 and S10 are not included in the 64 formal cells. Formal execution used no hard wall-time cap, and all 64 cells produced scored reports.

Let $y_{tar}$ be the original score for task $t$, arm $a$, and repeat $r$. We calculate $\bar y_{ta}=(y_{ta1}+y_{ta2})/2$, $d_t=\bar y_{t,\mathrm{skill}}-\bar y_{t,\mathrm{no}}$, and $\widehat\Delta=16^{-1}\sum_t d_t$. Thus tasks receive equal weight. The focal percentile bootstrap resamples paired tasks 10,000 times with mulberry32 seed 20260915. It quantifies variability over the observed tasks under resampling assumptions, not uncertainty over arbitrary repositories or model versions.

The report-judge-v2 rubric assigns a weight to each criterion. A pass receives its full weight, partial receives half, and fail or missing receives zero; declared triggered caps can limit the total. Criteria and frozen reference excerpts are supplied to the judge, which returns decisions and reasons rather than a total. The deterministic scorer calculates the 0--100 reward. Proposed verification can satisfy a static task; describing a test does not establish that it was executed.

The same model family supplies semantic evaluation. The run records describe arm-blind staging with one report per judgment; our later review was not blind. The archive retains original verdicts, including historical transport metadata; corrected provenance records distinguish harness-declared judge identity from independently observed provider identity. Scores are recomputed from rubric decisions. An integrity audit checks all 64 report hashes and all 328 original criterion decisions. Two stored reports require reversal of documented relative-link rewriting to reproduce their original hashes; this transformation is explicit rather than silently ignored.

\subsection{Bounded output review}
The initial targeted review was AI-assisted and non-blind, separate from the original verdicts. An earlier purposive pass covered five complete answers selected for prominent effects or operational concerns. A subsequent selection was committed before reading the additional full reports: paired answers for a positive task (S6), a negative task (S18), and a zero-change task (S12), at fixed repeats. One answer overlapped with the earlier pass. The combined scope is ten unique answers and 56 criteria.

The initial reviewer had access to task rubrics, reports, and earlier judgments, and selection used known outcomes. The authors subsequently reported human checking with broadly similar assessments; the later repository revision identifies contributing plugin authors as non-blind reviewers. The ten archived answers retain the same 56 decisions and replacement scores across these revisions. No separate per-reviewer scoring sheet or item-level confirmation was supplied, so revision stability is not interpreted as measured human--AI agreement or independent blinded validation. Both the initial AI-assisted analysis and the reported human follow-up are retained in the provenance ledger. Selection, report hashes, criterion-level reasons, and replacement scores are recorded. We keep original results as the primary recorded endpoint and show both the earlier and the combined replacement sensitivities. Replacements affect only reviewed answers; unreviewed scores are not presumed correct.

\subsection{Complete contract-stratified analysis}
We map each of the focal tasks' existing rubric criteria to one primary domain: version/release applicability; API/data/ownership; lifecycle/ordering/deployment; safety/boundary handling; failure attribution; or evidence/verification/provenance. The mapping is retrospective and fixed before generating the domain table. Multi-part criteria receive one primary label to avoid double counting, so these categories are an analytical organization of the rubric rather than an independently validated failure taxonomy.

Every formal report is linked to its task packet and original criterion verdicts through stored hashes. The analysis accounts for all 64 reports and all 328 criterion decisions, including original successes and shortfalls. We report full-credit counts with per-arm criterion denominators and distinct task counts. These are recorded judge outcomes, not newly verified correctness labels. This full-cohort accounting is distinct from the targeted ten-answer review; it does not represent an independent semantic regrading of all 64 answers. Repeats are nested within tasks, and tasks can contribute to multiple domains. We make no domain-level significance claims.

\section{Results}
\subsection{RQ1: Recorded reward changes}
\input{../generated/paired-effect-table.tex}
Historical mean changes range from $-3.07$ to $+10.67$ points (Table~\ref{tab:paired-effect}). The older Qwen configuration has substantial skill-arm timeouts under its local infrastructure. Terra excludes H8 because both arms have verifier timeouts. The GLM groups mix semantic and keyword judgments. These differences prohibit attributing cross-row variation to model capability alone. A contaminated Luna baseline, which accessed a native plugin-creator skill on H2 and H3, is excluded from this five-configuration comparison.

Within the historical GLM pair, the Flash-minus-5.2 difference in mean task lift is 6.23 points, with a paired bootstrap interval [2.59, 10.05]. The round-specific Flash gains are 6.32, 8.82, and 9.82, versus 1.55, 1.64, and 5.00 for GLM-5.2. Leaving one task out yields lift gaps from 5.33 to 7.00. These completed robustness calculations support a descriptive contrast within the recorded protocols. Different judge mixes and development histories remain confounders; the contrast does not identify a capability-dependent response curve.

\begin{table}[t]
\centering\small
\begin{tabular}{lrrrl}
\toprule
Focal endpoint & No skill & Skill & $\Delta$ & 95\% task-bootstrap interval\\
\midrule
Original judgments &93.83&98.75&4.92&[0.31, 10.86]\\
Earlier S11 replacement &93.83&98.44&4.61&[$-0.23$, 10.63]\\
All reviewed replacements &93.05&98.44&5.39&[0.00, 11.80]\\
\bottomrule
\end{tabular}
\caption{Focal 16-task comparison. Original scores remain unchanged in the archive. Replacement rows are partial-review sensitivity analyses, not validated estimates from a fully regraded dataset.}
\label{tab:focal}
\end{table}

In the focal comparison, mean original reward increases from 93.83 to 98.75 (Table~\ref{tab:focal}). Six tasks improve, two decline, and eight have zero change; all eight zero-change pairs are at 100 in both arms. S1 has a 42.5-point task-level gain. Removing it reduces the overall gain to 2.42 points across the remaining 15 tasks. The original signed-rank approximation gives $p=0.0797$ with eight nonzero pairs. The positive bootstrap lower endpoint should therefore not be portrayed as uniformly strong inferential support.

A supplementary newer Qwen S16 run reports 71.17 to 77.89, a gain of 6.72 points, with a reported interval [$-12.42$, 26.25]. It contains 64 scored attempts, including five timeouts and two other nonzero-exit attempts. Termination status is distinct from the assigned score. Its raw answers and grading reasons are not archived, and the reported bootstrap randomization is insufficiently specified for exact interval reproduction. Fourteen tasks overlap with the focal Flash run; different hosts, budgets, judges, and selections prevent a controlled cross-model comparison. We use this result only as supplementary context.

Three supplementary GLM-5.3 rounds on S1--S22 are now archived. Recomputing all 132 verdicts against the pinned rubric, retaining half-point values before summation, gives round gains of 1.93, 3.18, and 6.59 points. The per-task median across rounds is 97.05 without the skill versus 98.64 with it, a gain of 1.59 points; its descriptive task-bootstrap interval is [$-1.36$, 5.80] (10,000 samples, seed 20260907). The third round used GLM-5.3 as judge rather than the GLM-5.3-Flash judge used in rounds one and two. The median cannot remove this protocol change, so these results remain a mixed-judge descriptive supplement.

Round three's S17 no-skill answer receives zero because its affirmative immediate-registration recommendation triggers the rubric's contradiction cap. The criterion decisions are retained: this pair contributes 100 of the round's 145 total gain points. Excluding S17 descriptively leaves a 2.14-point mean gain across 21 tasks; this is a sensitivity, not a basis for removing the observation. The archive discloses quota-related relaunches, concurrency changes, and a truncated round-one answer. These rounds do not establish a capability ladder or restore comparability with older protocols.

\subsection{RQ2: Contract failures and grading disagreements}
% Generated by analyze-contract-coverage.mjs; do not edit.
\begin{table}[t]\centering\small
\begin{tabular}{p{0.37\textwidth}rrrr}
\toprule
Primary domain & Tasks & Criteria/arm & No skill & Skill\\\midrule
Version and release applicability & 6 & 12 & 8 & 12 \\
API, data and ownership contracts & 13 & 46 & 35 & 45 \\
Lifecycle, ordering and deployment & 9 & 26 & 24 & 25 \\
Safety and boundary handling & 5 & 10 & 10 & 10 \\
Failure attribution & 11 & 32 & 32 & 32 \\
Evidence, verification and provenance & 14 & 38 & 31 & 36 \\
\bottomrule\end{tabular}
\caption{Complete contract-stratified analysis of the focal 64 reports. The final two columns count original full-credit criterion decisions; they are not independently validated correctness counts. Criteria/arm includes both repeats. Multi-contract criteria receive one retrospective primary label; tasks can contribute to multiple domains.}
\label{tab:contracts}\end{table}

Table~\ref{tab:contracts} includes the full cohort. Original full-credit counts rise from 8 to 12 out of 12 version-related decisions and from 35 to 45 out of 46 API/data/ownership decisions per arm. Failure attribution has full credit in both arms (32/32), as does the safety/boundary domain (10/10). These counts describe the original judge; the reviewed S11 containment defect demonstrates why an all-full-credit domain cannot be called verified safe. Domains contain different tasks and weights, so their percentages are not a difficulty ranking or an estimate of a causal skill mechanism.

The bounded targeted review finds three criterion disagreements in two of ten answers. Because the sample is purposive and non-blind, these counts are not an error-rate estimate. They identify concrete weaknesses in the recorded measurement.

\paragraph{Information access and public API boundaries.}
S1's reviewed answers retain their original 55 and 100 scores. Part of its rubric rewards mapping to supplied migration cards. The large gain therefore includes access to organized facts and identifiers, and cannot isolate a procedural reasoning benefit. In S6, the no-skill answer correctly recommends removing obsolete defensive code but incorrectly attributes informational-event handling wholly to the host. It also recognizes a public append API gap while recommending an unqualified append route. Review changes two criteria and lowers this answer from 87.5 to 62.5. The paired skill answer retains 100 by distinguishing producer responsibilities, retention, and unavailable public surfaces. The resulting sensitivity increases estimated skill benefit, illustrating that grading errors can favor either arm. Additional full-text inspection of S6 repeat 2 finds the same platform-only interpretation in the no-skill answer, whereas both skill answers distinguish retention from a supported producer seam. This is repetition within one task, not a second independent incident; no additional score replacement is imposed in the fixed ten-answer sensitivity.

\paragraph{Containment boundaries.}
In S11, a with-skill answer proposes a relative-path check equivalent to accepting an empty relative path or any non-absolute path that does not start with two dots followed by a separator. The exact relative path \texttt{..} passes this predicate, even though it denotes the parent directory. An offline counterexample exercises this case under both POSIX and Windows path semantics. Review lowers the containment criterion from full to partial credit and the answer from 90 to 80. This establishes a defect in the proposed guard; it does not establish a complete exploitable endpoint or claim that every deployment has the same exposure.

\paragraph{Cleanup versus process liveness.}
For S18 repeat 2, the no-skill answer retains 100 and the skill answer retains 90. The latter supplies component-unmount cleanup but does not satisfy the timer-liveness requirement to avoid keeping the process alive. Cleanup and timer unreferencing address different conditions. The original partial credit is retained. This example demonstrates why a plausible lifecycle recommendation is not interchangeable with the requested runtime contract.

\paragraph{Ceiling and a retained partial score.}
Both reviewed S12 answers retain 100: they connect a Windows file lock to the native host process, distinguish the recorded release channels, and prescribe the pinned installation and restart sequence. This zero-change pair illustrates a ceiling case rather than evidence that the skill was consumed or unnecessary in all contexts. The previously reviewed S17 no-skill answer retains 80 because recognizing onboarding requirements does not repair an incorrect injection descriptor. No frequency or causal mechanism is inferred from these selected examples.

\subsection{RQ3: Sensitivity and resource use}
Replacing only the reviewed S11 score changes the focal estimate to 4.61 points and moves its interval across zero. Replacing all reviewed scores changes it to 5.39 points with a lower endpoint at zero. The direction remains positive, but the strength of evidence depends on grading decisions. Neither the larger replacement estimate nor the original positive interval should be selected in isolation.

The 64 formal execution records contain 3,185,993 no-skill and 12,320,379 skill \texttt{subagent\_tokens}, a ratio of 3.87. The corresponding sums of recorded task durations are 9,446 and 11,450 seconds, a ratio of 1.21. The token field does not separately define input, output, and cache accounting, so it is not a billing-cost measure. Summed task durations are not end-to-end wall-clock time under overlapping execution. Pilot records are excluded.

We also recomputed the earlier Flash round-one usage totals from its 22 sessions per condition. Fresh input totals are 599,882 and 678,101; output totals are 116,506 and 132,726; cache totals are 5,890,176 and 5,017,984. The recorded total-token sums are 6,606,564 and 5,828,811, respectively. These fields do not support an undifferentiated claim of a 2.2-fold token increase. Their accounting differs from the focal host's token field, so resource ratios are reported within configurations only.

\section{Discussion}
\subsection{What the evidence supports}
The archive supports a bounded claim: this version-pinned skill is associated with higher mean recorded reward in several evaluated configurations, while benefits are uneven and selected answers still violate operational contracts. In the focal comparison, ceiling effects and one large task gain limit the interpretation of its average. The review adds practical detail by showing where apparently adequate explanations fail at the boundary between migration knowledge and an actionable repair.

For maintainers, the relevant implication is to verify the target contract explicitly: whether an API is publicly available, whether a containment predicate covers the parent itself, and whether a timer permits process exit. A correct vocabulary match or a plausible cleanup pattern is insufficient. These are recommendations grounded in the observed examples, not a measured intervention showing that adding these checks improves future agents.

\subsection{What the evidence does not support}
The historical ordering is not evidence of an inverted U. Baseline reward is an outcome on the same tasks, shares mathematical components with the gain, and has a bounded ceiling. It is not an independently measured capability axis. Cross-model differences additionally mix infrastructure, graders, task exposure, and skill versions. An inverted-U claim would require a separate design; it is not a prerequisite for the present retrospective contribution.

Likewise, this study does not isolate the effect of procedural guidance from additional information, estimate production repair reliability, establish cost effectiveness, or show generalization to unseen projects. Prior skill benchmarks already address broader effectiveness questions. Our contribution depends on transparent linkage among migration contracts, archived answers, scoring decisions, and qualified conclusions in one setting.

\section{Threats to Validity}
\paragraph{Construct validity.}
Static diagnostic reward is not functional repair success. Some criteria reward card mappings or metadata. Keyword and semantic graders measure different things, and numeric normalization does not harmonize their constructs. The focal solver and judge share a model family. The initial non-blind AI-assisted review is useful for finding counterexamples but does not certify all scores. Human follow-up is author-reported without a separate item-level scoring dataset, precluding an inter-rater agreement estimate. Full-cohort contract tabulation inherits the original judge decisions and should not be confused with complete revalidation.

\paragraph{Internal validity.}
The data are retrospective, and contributors developed the skill and tasks in an overlapping workflow. Historical selections are not assumed random, and the focal seeded draw does not remove development exposure. Historical execution budgets, environment failures, and protocols differ. Original model labels rely on recorded infrastructure declarations. The review selection was frozen before additional full-output inspection, but after outcomes were known; it is not preregistration of the underlying experiment.

\paragraph{Statistical conclusion validity.}
There are only 16 focal tasks and two repeats per arm. Eight task pairs are at the ceiling. Task-level resampling treats tasks as independent even when incident sources may overlap. The selected review does not estimate population error, and partial replacements leave unreviewed scoring uncertainty unresolved. Intervals, signed-rank approximations, and leave-one-out results are exploratory descriptions of the recorded data rather than confirmatory tests of a universal effect.

\paragraph{External validity and reproducibility.}
The focal tasks concern one plugin framework and static reports. Historical task inventories and artifacts have unequal coverage. The newer Qwen outputs and grading reasons are unavailable in the repository, limiting replication of its endpoint. Live models and dependencies can change even with pinned task sources. The archive supports deterministic recalculation of the focal scores and reported sensitivities, not a guarantee of identical future model outputs.

\section{Conclusion}
A migration skill can improve recorded diagnostic reward while leaving important operational requirements unmet. In a completely archived 16-task comparison, the observed gain is 4.92 points, but ceiling effects, concentration in one task, and grading sensitivity qualify its interpretation. Complete contract-stratified accounting locates differences in recorded credit, while a bounded targeted review finds disagreements in both conditions and exposes distinctions between available APIs and suggested actions, plausible containment checks and boundary correctness, and cleanup and process liveness. The resulting contribution is a traceable retrospective maintenance study, with explicit limits on grading validity and generalization. It neither requires nor establishes an inverted-U capability relationship.

\section*{Artifact Availability and Research Transparency}
The working repository is \url{https://github.com/oh-my-dsh/dsh-plugin-upgrade-skill}. The focal archive is under \path{benchmark/results/artifacts/2026-09-15-glm-5.3-flash-unified-s16}; the supplementary review is under \path{paper/audit/output-review-20260917}. Original outputs and scores are preserved separately from review verdicts. Complete contract accounting and human-follow-up provenance are stored under \path{paper/audit/contract-review-20260917}. The claim-to-evidence ledger identifies inputs and offline commands for regenerating the analyses. Availability is described per configuration; no universal raw-output completeness claim is made. A permanent release identifier and redistribution review remain author submission checks.

\section*{AI Assistance and Author Declarations}
AI tools (Codex and ChatGPT) assisted with source inspection, analysis code, initial targeted output review, and manuscript drafting. The authors subsequently reported non-blind human checking by contributing plugin authors. The archive preserves both stages and does not claim independently measured human agreement.

Author identities, contributions, approval, funding, competing interests, and the final publisher-specific AI declaration must be confirmed by the authors before submission; this review draft does not assert those confirmations.

\bibliographystyle{plainnat}
\bibliography{custom}
\appendix
\clearpage
\section{Task-level focal results}
% AUTO-GENERATED from the original 64 scored cells; do not edit.
\begin{longtable}{lrrrrr}
\caption{Focal task-level original scores. Each arm has two repeats; differences use arm means. These are static rubric scores, not repair pass rates.}\label{tab:focal_tasks}\\
\toprule
Task & No skill r1 & No skill r2 & Skill r1 & Skill r2 & Mean $\Delta$\\
\midrule\endfirsthead
\toprule Task & No skill r1 & No skill r2 & Skill r1 & Skill r2 & Mean $\Delta$\\\midrule\endhead
S1 & 55.00 & 60.00 & 100.00 & 100.00 & 42.50 \\
S2 & 100.00 & 80.00 & 100.00 & 100.00 & 10.00 \\
S5 & 87.50 & 100.00 & 100.00 & 100.00 & 6.25 \\
S6 & 87.50 & 87.50 & 100.00 & 100.00 & 12.50 \\
S9 & 100.00 & 100.00 & 100.00 & 100.00 & 0.00 \\
S11 & 100.00 & 90.00 & 90.00 & 90.00 & -5.00 \\
S12 & 100.00 & 100.00 & 100.00 & 100.00 & 0.00 \\
S13 & 100.00 & 100.00 & 100.00 & 100.00 & 0.00 \\
S14 & 100.00 & 100.00 & 100.00 & 100.00 & 0.00 \\
S15 & 100.00 & 100.00 & 100.00 & 100.00 & 0.00 \\
S17 & 80.00 & 80.00 & 90.00 & 100.00 & 15.00 \\
S18 & 100.00 & 100.00 & 100.00 & 90.00 & -5.00 \\
S19 & 100.00 & 100.00 & 100.00 & 100.00 & 0.00 \\
S20 & 97.50 & 97.50 & 100.00 & 100.00 & 2.50 \\
S21 & 100.00 & 100.00 & 100.00 & 100.00 & 0.00 \\
S22 & 100.00 & 100.00 & 100.00 & 100.00 & 0.00 \\
\bottomrule
\end{longtable}


\section{Artifact map and reproduction scope}
Historical result files are stored under \path{benchmark/results}: \path{paired-effect-stats.json} generates the paired table, and \path{glm-pair-stability.json} records GLM stability. The focal directory contains the selection, schedule, execution-order records, original reports and judgments, aggregate scores, paired analysis, and the earlier targeted review. The later review stores frozen selection and complete criterion verdicts separately. The script \path{paper/scripts/summarize-submission-evidence.mjs} checks report hashes, recomputes review scores, deduplicates overlapping answers, and regenerates sensitivity and historical resource summaries without model calls. Its \texttt{--check} mode verifies the committed generated summary.

A frozen 23-task historical inventory and its source metadata are retained in \path{paper/generated}; they are not substituted for the 16-task selection. Exposure information remains in \path{paper/audit/task-exposure-ledger.csv}. Missing raw artifacts, unverified model identities, and pending author declarations are recorded as limitations rather than filled in by inference.
\end{document}
